\noindent
This dissertation examines farmer storage behavior in agricultural procurement markets characterized by localized, time-varying buyer power. Combining institutional analysis, theoretical modeling, and micro-level field data from China's fresh apple sector, the analysis demonstrates that when buyer power is dynamic rather than fixed, farmers may benefit from intertemporal storage and marketing strategies that allow access to more competitive selling conditions.

\section{Summary of Main Findings}
\noindent
The descriptive chapter documents three institutional features of China's apple procurement markets. First, apple marketing is highly localized, with effective buyer competition occurring at the village or county level rather than across broader regions. Second, despite the presence of many nominal buyers during harvest, trading activity is frequently characterized by coordination and collusion among intermediaries. Third, storage is widespread among farmers but highly heterogeneous in efficiency, access costs, and scale.

The theoretical chapter reframes oligopsony power as a time-varying phenomenon. Rather than assuming buyer power is constant, the model allows buyer competitiveness to evolve across selling periods as traders enter, exit, or reallocate procurement activity. Within this framework, storage generates an option value: farmers with relatively efficient storage may gain access to future selling windows with improved competition. The analysis shows that temporal variation in buyer power alone can induce storage, even absent predictable seasonal price movements. These gains, however, are substantially reduced by storage inefficiency and farmer risk aversion.

The empirical chapter provides field-level evidence consistent with the model's predictions. Using household survey data from Fuji apple growers, the analysis shows that storage decisions respond primarily to perceived buyer competitiveness rather than to the observed number of buyers at harvest. Expectations of increased future competition significantly raise storage adoption and delay sales, while expectations of fewer buyers have little deterrent effect. This asymmetric response highlights the importance of anticipated market opportunities in shaping farmer behavior.

Taken together, the findings indicate that farmer marketing decisions reflect forward-looking responses to evolving buyer conditions, rather than passive reactions to contemporaneous prices or buyer counts.

\section{Policy Implications}
\noindent
Public investment in agricultural storage is commonly justified by objectives such as reducing post-harvest losses or smoothing seasonal price fluctuations. The results here identify an additional channel: storage alters farmers' exposure to buyer power by enabling intertemporal access to potentially more competitive selling environments.

However, storage does not uniformly improve farmer welfare. Storage generates economic value only when efficiency losses are limited and when farmers can bear associated price and market uncertainty. Expanding storage capacity without addressing efficiency, access costs, or risk exposure may therefore yield limited income gains. Policies that improve storage efficiency, reduce transaction and access costs---particularly distance-related barriers---and support shared-use storage arrangements are more likely to enhance storage effectiveness.

A central empirical finding is that farmers respond to perceived and expected buyer competitiveness rather than to observed buyer counts. This distinction has important implications for market information policy. In environments characterized by trader coordination or collusion, structural indicators such as buyer numbers may substantially overstate effective competition. Farmers instead rely on experience, informal information, and social networks to form expectations about future market conditions.

Improving market transparency may therefore influence behavior primarily by shaping expectations rather than by directly affecting prices. Timely dissemination of price information, buyer participation patterns, and historical seasonal trading conditions can reduce informational asymmetries and improve farmers' assessment of competitive opportunities.

The results also underscore the role of risk exposure in shaping storage behavior. Even when storage is technologically feasible, greater risk aversion substantially suppresses storage adoption and shortens marketing horizons. Complementary policies, such as improved access to credit, risk-sharing mechanisms, or price-risk management instruments, may therefore strengthen the efficiency of storage and enable farmers to tolerate uncertainty associated with delayed sales better.

Storage adoption further exhibits substantial heterogeneity across households. Education, leadership status, and social capital are positively associated with storage use, indicating that informational and institutional capacity conditions farmers' ability to engage in forward-looking marketing strategies. Storage access is thus shaped not only by physical infrastructure and financial resources, but also by information and networks.

Finally, the findings have direct implications for competition policy in agricultural procurement markets. The presence of multiple buyers does not guarantee competitive outcomes when coordination among traders is prevalent. Policies focused solely on increasing buyer numbers may therefore have limited effects. By contrast, mechanisms that expand farmers' ability to choose the timing of sales can indirectly discipline buyer power by strengthening farmers' outside options. In this sense, storage does not substitute for competition policy but interacts with it.